\section{Conclusion}

    Au cours de ce stage, j'ai exploré la littérature scientifique portant sur le RAG et son explicabilité. J'ai découvert les différentes perspectives d'évolution de ces systèmes et les perspectives d'explicabilité pour les tâches l'extraction d'information à l'aide des méthodes d'interprétabilité mécaniste. J'ai repris le système RAG proposé par \cite{tran_retrieval_2024} et conçu un dataset type à partir des données fournies par le challenge RAG@EvalLLM. Enfin, j'ai proposé une méthodologie d'usage de l'activation patching avec l'outil MechIR \cite{parry_mechir_2025}. Cette dernière inclus la conception des perturbations ainsi que la visualisation et l'interprétation des résultats fournis par MechIR.

\section{Perspectives générales}

    Au cours du stage, j'ai finalement choisi de poursuivre un sujet de thèse différent de celui visé par ce stage. Les travaux réalisés  constituent néanmoins une première exploration méthodologique autour de l'interprétabilité mécaniste appliquée aux modèles d'extraction d'information pour le RAG.

    Dans l'éventualité d'une reprise de ces travaux, voici quelques perspectives pour tende vers le résultat initialement prévu. Dans un premier temps, L'identification des nœuds proposée ne constitue pas encore une cartographie d'un modèle d'extraction. C'est la première étape. Une fois la cartographie obtenue, l'extraction d'explication du fonctionnement du modèle lors de l'usage doit être définie. Il est nécessaire de garder en tête que l'utilisateur visé n'est pas informaticien. Par conséquent, un travail doit être effectué pour rendre les explications accessible au public visé.

    Par ailleurs, d'autres perspectives d'usage d'une telle cartographie existent. Cette dernière peut servir d'orientation pour améliorer ou spécialiser un modèle. De premiers résultats existent déjà sur l'utilisation de l'activation patching appliqué au prunning.